\subsection{Кинематика  колебаний}
%%%%%%%%%%%%%%%%%%%%%%%%%%%%%%%%%%%%%%%%%%%%%%%%%%%%
%%%%%%%%%%%%%%%%%%%%%%%%%%%%%%%%%%%%%%%%%%%%%%%%%%%%
 \z{Груз, свободно колеблющийся на пружине, за время $0.01$ с сместился с расстояния $0.5$ см от положения равновесия до наибольшего, равного $1$ см. Каков период его колебаний?}
%
{$T = 0.06$ c.}
%%%%%%%%%%%%%%%%%%%%%%%%%%%%%%%%%%%%%%%%%%%%%%%%%%%% 
 \z[{\begin{wrapfigure}[8]{r}{0.2\textwidth}%
		\vspace{-10pt}%
		\begin{flushright}%
 			\includegraphics[width=0.2\textwidth]{10.3.2}%
		\end{flushright}%
	\end{wrapfigure}%
	}]%
	{К наклонной стене подвешен маятник длины $l$. Маятник отклонили от вертикали на малый угол, в два раза превышающий угол наклона стены к вертикали, и отпустили. Найдите период колебаний маятника, если удары о стену абсолютно упругие.}
%
{$T = \frac{4}{3}\pi\sqrt{\frac{l}{g}}$}
%%%%%%%%%%%%%%%%%%%%%%%%%%%%%%%%%%%%%%%%%%%%%%%%%%%% 
 \z{Точечная частица совершает гармонические колебания вдоль оси $Ox$. В некоторый момент времени частица имеет координату $x_1 = 3.0$ см, скорость $v_{x1} = 8.0$ см/с и ускорение $a_{x1} = -12$ см/с$^2$. Определите амплитуду $A$, циклическую частоту $\omega$ и период $T$ колебаний. Начало координат совпадает с положением равновесия частицы.}
%
{$A = 5$ см, $\omega = 2$ рад/с, $T = 3.14$ c.}
%%%%%%%%%%%%%%%%%%%%%%%%%%%%%%%%%%%%%%%%%%%%%%%%%%%%	
 \z{Частота свободных колебаний тела равна $\omega$. Через какое наименьшее время его кинетическая энергия уменьшается вдвое по сравнению со своим наибольшим значением?}
%
{$t = \frac{\pi}{4\omega}$}
%%%%%%%%%%%%%%%%%%%%%%%%%%%%%%%%%%%%%%%%%%%%%%%%%%%%
\z{Точечная частица совершает гармонические колебания вдоль оси $Ox$, с циклической частотой $\omega = 4.0$ с$^{-1}$. В некоторый момент времени частица имеет координату $x_1 = 25$ см и скорость $v_{x1} = 1.0$ м/с. Определите координату $x_2$ и скорость $v_{x2}$ частицы спустя $\Delta t = 2.4$ с. Начало координат совпадает с положением равновесия частицы.}
%
{$x_2 = -28.9$ см, $v_{x2} = 81$ см/с.}